\section{The concentration and arithmetic Gram estimates}
\label{sec:cutoff-estimates}
This section establishes the cutoff estimates needed for the finite-packet
argument. All conductors are fixed before the cutoff tends to infinity.
The raw parity space is identified with $L^2((0,\infty),du)$ by even
or odd extension, with the positive-half-line norm. The raw Fourier
transform is the restriction of the transform with phase
$e^{-2\pi iu\xi}$ to that parity space; it is unitary in this norm.
Let $P_\Lambda$ be raw time restriction and let $B_R$ be raw band
restriction after Fourier transformation, where $R=\Lambda/q$.
Write $c=\Lambda R$.

For raw functions put $(Uf)(t)=e^{t/2}f(e^t)$ and
$\mathcal Mf(s)=\int_0^\infty f(u)u^{s-1}du$ wherever the integral
converges. Then $\mathcal Mf(1/2+i\tau)=\widehat{Uf}(\tau)$.
The map $U$ is unitary to $L^2(\R,dt)$.
We use $H_D^s=U^{-1}H^s(\R)$, where
\[
 H^s(\R)=\left\{g:\int_\R(1+\tau^2)^s
                     |\widehat g(\tau)|^2\frac{d\tau}{2\pi}<\infty\right\},
\]
with the norm given by this integral's square root. For negative $s$
the definition is understood in tempered distributions; only
nonnegative $s$ are needed below.

The projection $\Pi_\Lambda$ in this section always means the
\emph{unconstrained} parity selection of the time-supported eigenvectors
of $P_\Lambda B_RP_\Lambda$ with full prolate indices $n\leq4c$.
In the principal sector the raw column space used for the theta image
is instead $\mathcal R_\Lambda^0=\operatorname{Ran}\Pi_\Lambda
\cap\ker I\cap\ker\operatorname{ev}_0$.
Projection onto this smaller space is denoted $\Pi_\Lambda^0$.
The distinction is necessary: the interior prolate approximation
applies first to $\Pi_\Lambda$, before the two constraints are restored.

\begin{lemma}\label{lem:raw-estimates}
For sufficiently large $\Lambda$, the selected concentration eigenvalues
are at least $1/4$, and the first omitted eigenvalue is at most $3/4$.
There are $d>0$ and constants $C_s$, independent of $\Lambda$, such that
every $\psi\in\operatorname{Ran}\Pi_\Lambda$ satisfies
\[
 \int_{d/R}^\Lambda|\psi(u)|^2du\geq\tfrac12\|\psi\|_2^2,
 \qquad \|U\psi\|_{H^s}\leq C_s(1+c)^s\|\psi\|_2
 \quad(0<s<1/2),
\]
where $(U\psi)(t)=e^{t/2}\psi(e^t)$ and all time-supported functions
are extended by zero. The conclusions pass to the constrained space.
\end{lemma}
\begin{proof}
Landau's transition formula at angular bandwidth $\pi$ is
\[
 \nu_{\lfloor2T+b\log T\rfloor}\longrightarrow
                    (1+e^{\pi^2b})^{-1}.
\]
See \citet[pp.~209, 211]{landau1985}; the underlying distribution
theorem is \citet{landauwidom1980}. Scaling gives $T=2c$.
Monotonicity and a squeeze with $b=\pm\varepsilon$ show that the
eigenvalue at $\lfloor4c\rfloor$, and at each bounded offset from it,
tends to $1/2$. This includes the next index of the selected parity.
The stated bounds follow.

Put $A=P_\Lambda B_RP_\Lambda$ on the selected finite space and
$w=B_RA^{-1}\psi$. Then $P_\Lambda w=\psi$ and
$\|w\|_2^2=\langle\psi,A^{-1}\psi\rangle\leq4\|\psi\|_2^2$.
Band support, Cauchy--Schwarz in Fourier inversion, and Plancherel
give $|w(u)|^2\leq C R\|\psi\|_2^2$ and
$\|w'\|_2\leq C R\|\psi\|_2$.
Integrating the first bound over $(0,d/R)$ and choosing $d$ sufficiently
small proves the mass assertion.

Let $f=U\psi$, considered first on $(-\infty,L)$,
$L=\log\Lambda$. Its interior derivative has norm
\[
 \|f'\|_{L^2(-\infty,L)}
 \leq\tfrac12\|\psi\|_2+\Lambda\|w'\|_2
 \leq C(1+c)\|\psi\|_2.
\]
Its value tends to zero at the lower end. The fundamental theorem
of calculus and Cauchy--Schwarz imply
$|f(L-)|^2\leq2\|f\|_2\|f'\|_2$.
The Fourier transform of the zero extension obeys the distributional
identity $-i\tau\widehat f(\tau)=\widehat{f'1_{t<L}}(\tau)
-f(L-)e^{i\tau L}$. For $M\geq1$, its contribution on $|\tau|>M$
to the squared $H^s$ norm is therefore at most
\[
 C_s\left(M^{2s-2}\|f'\|_2^2+
                    M^{2s-1}|f(L-)|^2\right).
\]
The boundary integral converges because $s<1/2$.
The contribution on $|\tau|\leq M$ is at most
$C M^{2s}\|f\|_2^2$. Taking $M=1+c$ proves the Sobolev estimate,
including the sharp endpoint jump.
\end{proof}

\begin{lemma}\label{lem:theta-identities}
When a positive argument $x$ is used, $E_\chi\psi(x)$ denotes
the logarithmic function $(E_\chi\psi)(\log x)$.
On smooth compact raw functions of the required parity, the centered
theta map has Mellin multiplier $L(1/2+i\tau,\chi)$ and
\[
 \|E_\chi\psi\|_{H^{1/16}(dt)}
      \leq C_\chi\|U\psi\|_{H^{3/8}(dt)}.
\]
For nonprincipal primitive $\chi$, with Gauss sum $\tau(\chi)$,
\[
 E_\chi\psi(x)=\frac{\tau(\chi)}{\sqrt q}
     E_{\bar\chi}\widehat\psi\bigl((qx)^{-1}\bigr).
\]
For the principal character the corrected map
$E_1\psi(x)=\sqrt x\sum_{n\geq1}\psi(nx)-x^{-1/2}I\psi$
satisfies $E_1\psi(x)=E_1\widehat\psi(x^{-1})$.
These identities extend to the graph domains and the finite selected
spaces used below.
\end{lemma}
\begin{proof}
For the uncorrected theta sum in an absolute-convergence half-plane,
the change of variables $u=nx$ in the Mellin integral gives
$L(s,\chi)\mathcal M\psi(s)$.
For nonprincipal $\chi$, bounded character partial sums and Abel
summation continue this identity to $\Re s>0$ where the integral
converges. In the principal case split the uncorrected Mellin
integral at $x=1$, subtract its mean term on $(0,1)$, and continue
that term as $I\psi/(s-1)$. The upper-tail integral of the
corrected map gives the same expression on $0<\Re s<1$.
Euler summation justifies the remaining integral there. A bounded
variation compact raw function has corrected theta sum $O(\sqrt x)$
as $x\downarrow0$; its upper tail is zero in a nonprincipal sector
and is $-I\psi\,x^{-1/2}$ in the principal sector. Differentiated
Mellin integrals therefore converge on $0<\Re s<1$.

The Euler product on $\Re s=1+\varepsilon$, the functional equation,
Stirling's formula, and the Phragm\'en--Lindel\"of principle in the
strip $-\varepsilon\leq\Re s\leq1+\varepsilon$ give
$|L(1/2+i\tau,\chi)|\leq C_\chi(2+|\tau|)^{5/16}$.
For example use $\varepsilon=1/8$; the left boundary has exponent
$1/2+\varepsilon$ and interpolation at the midpoint halves it.
For zeta first cancel its pole by a rational factor bounded above
and below on the critical line. Bounded heights are absorbed in
the constant. Multiplication by this bound and Plancherel give the
displayed estimate, since $5/16+1/16=3/8$.

For Poisson summation, expand the periodic character on residue
classes modulo $q$. Its finite Fourier coefficients are
$\tau(\chi)\overline\chi(m)$. Applying the ordinary Poisson formula
on each class yields
\[
 \sum_{n\in\Z}\chi(n)\psi(nx)=
    \frac{\tau(\chi)}{qx}\sum_{m\in\Z}
                \overline\chi(m)\widehat\psi(m/(qx)).
\]
The parity condition makes each side twice its positive-index sum.
Multiplication by $\sqrt x$ gives the assertion.
In the principal even case the zero terms in the ordinary Poisson
formula are $\psi(0)$ and $\widehat\psi(0)=2I\psi$;
also $I\widehat\psi=\psi(0)/2$. Substitution proves the corrected
identity. Smooth approximation in $H_D^{5/16}$, followed by the
multiplier estimate, extends the identities whenever the specified
support or finite constraints are preserved. The raw Fourier
transform reverses the self-adjoint dilation generator and thus
preserves $H_D^s$, proving graph convergence on both sides. On a band space,
evaluation at zero is continuous by Fourier Cauchy--Schwarz.
\end{proof}

\begin{proposition}[Unconditional arithmetic Gram control]
\label{prop:gram-control}
Let $G_\Lambda$ be the Gram matrix of the theta images of a raw
orthonormal basis of the selected nonprincipal space or of
$\mathcal R_\Lambda^0$ in the principal sector. For every $B>0$,
\[
 \|G_\Lambda^{-1}\|\leq C_{\chi,B}
                          c(\log c)^{-2B}
\]
for sufficiently large $\Lambda$. No zero-location theorem is used.
\end{proposition}
\begin{proof}
The arithmetic input is the fixed-character specialization of
\citet[Lemma 3.9(ii), pp.~46--47]{mstt2023}:
for every fixed $K,B>0$,
\[
 \sup_{|\tau|\leq N^K}
 \left|\sum_{n\leq N}\frac{\mu(n)\chi(n)}{n^{1/2+i\tau}}\right|
        \leq C_{\chi,K,B}\frac{\sqrt N}{(\log N)^B}.
\]
To obtain the displayed form from their interval version, discard
$n\leq2\sqrt N$ at cost $O(N^{1/4})$, partition the rest into
dyadic intervals $[Y,2Y]$, and apply the source with $X=Y^{2K}$
and exponent $\alpha=1/(2K)$, after increasing $K$ to at least one.
The square roots of the interval lengths sum to $O(\sqrt N)$;
the fixed modulus eventually satisfies the source's logarithmic
modulus condition. On the entire frequency line the trivial bound
for this polynomial is $2\sqrt N$.

Put $a=d/R$, $N=\lceil\Lambda/a\rceil\asymp c$,
$f=U\psi$, and $e=E_\chi\psi$. Since $I\psi=0$ in the
principal space, $e(t)=0$ for $t>L$ in all sectors. Finite M\"obius
inversion gives, exactly for $\log a\leq t\leq L$,
\[
 f(t)=\sum_{n\leq N}\frac{\mu(n)\chi(n)}{\sqrt n}
                                           e(t+\log n).
\]
Indeed expansion of the finite theta sums assigns to the term with
index $k$ the coefficient $\chi(k)\sum_{n\mid k}\mu(n)$.
All terms beyond the raw support vanish. This argument also holds
in local $L^2$, by approximation, without assigning point values
to an arbitrary $L^2$ function.

Let $\mathcal B_N$ denote the finite translation operator on the
right. Split its Fourier multiplier at $|\tau|=N^{16}$.
The arithmetic bound and the global trivial bound imply
\[
 \|\mathcal B_N e\|_2\leq
 C_B\frac{\sqrt N}{(\log N)^B}\|e\|_2
       +2\sqrt N\,N^{-1}\|e\|_{H^{1/16}}.
\]
The preceding two lemmas bound the last term by
$C c^{-1/8}\|\psi\|_2$, since $1/2-1+3/8=-1/8$.
The mass estimate and the exact inverse identity bound the left
side below by $2^{-1/2}\|\psi\|_2$. Absorb the last term for
large $c$, and square. The resulting inequality
\[
 \|E_\chi\psi\|_2^2\geq C_B^{-1}
                 c^{-1}(\log c)^{2B}\|\psi\|_2^2
\]
is exactly the claimed lower bound for the Gram matrix.
\end{proof}
